% !TEX root = main.tex
\section{Source Value Set}\label{sec:source-value-set}

Now let's from the definition of source value, derive two kinds of sets.
\begin{definition}~\label{def:source-value-set}
Given a basic target value set $B \subset \widehat{\mathbb{N}}_{{\beta}}$. Denote that
\begin{align*}
    \mathbb{S}_n(B) &= \bigcup_{b\in B}\left\{ f_n(\bm{v};b) \mid \bm{v} \in \mathbb{V}_n(b) \right\},\\
    \mathbb{V}_n(B) &= \bigcup_{b\in B}\left\{ \bm{v} \in \mathbb{N}^{n} \mid f_{n}(\bm{v};b) \in \mathbb{N} \right\}.
\end{align*}
Specially, we define $\mathbb{V}_0(B) = \varnothing$ and $\mathbb{S}_0(B) = B$. When $B = \{b\}$ contains only one number, we also denote as $\mathbb{V}_n(b)$ and $\mathbb{S}_n(b)$.
\end{definition}
\begin{proposition}
    $\mathbb{S}(B) = \bigcup_{n = 0}^{\infty}\mathbb{S}_{n}(B)$
\end{proposition}
\begin{proposition}~\label{prop:source-set-include}
    If $a\overset{n}{\rightsquigarrow} b$, $\mathbb{S}(a) \subset \mathbb{S}(b)$ and $\mathbb{S}_{k}(a) \subset \mathbb{S}_{k+n}(b)$.
\end{proposition}
\begin{proposition}~\label{prop:source-set-disjoint}
    If $a \not\rightsquigarrow b$ and $b \not\rightsquigarrow a$, $\mathbb{S}(a) \cap \mathbb{S}(b) = \varnothing$.
\end{proposition}
\begin{proposition}~\label{prop:non-cycle-layer-independent}
    If $b \not\rightsquigarrow b$, then $\forall n \ge 1$, $\mathbb{S}_{n}(b) \cap \bigcup_{k = 0}^{n-1}\mathbb{S}_{k}(b) = \varnothing$.
\end{proposition}
\begin{proof}
    If $\exists m \le n$, $\mathbb{S}_{m}(b) \cap \mathbb{S}_{n}(b) \ne \varnothing$, $\exists a \in \mathbb{S}_{m}(b) \cap \mathbb{S}_{n}(b)$ such that $a \rightsquigarrow a$. Notice that $a \rightsquigarrow b$, thus $b\rightsquigarrow b$. That's a confliction with $b \not\rightsquigarrow b$.
\end{proof}
\begin{proposition}~\label{prop:cycle-layer-included}
If $b \rightsquigarrow b$ in a cycle that period $t$, then $\forall k\ge 0$, $b \in \mathbb{S}_{kt}(b)$, and hence $\mathbb{S}_{kt}(b)\subset \mathbb{S}_{(k+1)\cdot t}(b)$.
\end{proposition}
From \cref{cor:modulo-alpha-of-f_1}, we still have
\begin{proposition}
    If $\alpha \mid \gamma$, then $\forall b\not\in\mathbb{S}_{\infty}$, we have $\mathbb{S}_{1}(b) \not\subset \mathbb{S}_{\infty}$.
\end{proposition}
\begin{proposition}~\label{prop:gcd-1-derive-not-all-primitive}
    Suppose $\gcd(\alpha, \nabla_1) = 1$. Then $\forall b\not\in\mathbb{S}_{\infty}$, $\forall n\ge 1$, we have $\mathbb{S}_{n}(b)/\alpha = \mathbb{Z}/\alpha$. And hence $\mathbb{S}_{1}(b) \not\subset \mathbb{S}_{\infty}$ and $\mathbb{S}_{1}(b) \cap \mathbb{S}_{\infty} \ne\varnothing$.
\end{proposition}
\begin{proof}
    We prove $n=1$ at first.\newline
    Suppose $a \in\mathbb{S}_{1}(b) \cap \mathbb{S}$, and its reduce vector is $\bm{a}$.
    From \cref{cor:modulo-alpha-of-f_1}, $\forall k\in\mathbb{Z}$ that $f_{1}(\bm{a} + k\Phi_1\bm{e}_1) > 0$, there is $f_{1}(\bm{a} + k\Phi_1\bm{e}_1) \equiv k\nabla_1 + f_{1}(\bm{a}) \pmod{\alpha}$ runs all remainders modulo $\alpha$. Thus $\mathbb{S}_{1}(b)/\alpha = \mathbb{Z}/\alpha$.\newline
    Now consider any layer $n > 1$. Since any $n$-step source value is just a $1$-step source value of a $n-1$-step source value. Thus they follow the conclusion of $n=1$. Thus we proved $\mathbb{S}_{n}(b)/\alpha = \mathbb{Z}/\alpha$.
\end{proof}
\begin{proposition}~\label{prop:all-primitive-derive-gcd-ge-1}
    If $\exists b\not\in\mathbb{S}_{\infty}$ such that $\mathbb{S}_{1}(b) \subset \mathbb{S}_{\infty}$, then $\gcd(\alpha, \nabla_1) > 1$.
\end{proposition}
\begin{proposition}~\label{prop:have-primitive-derive-all-primitive}
    Suppose $\alpha\mid \nabla_1$. If $\exists b\not\in\mathbb{S}_{\infty}$ such that $\mathbb{S}_{1}(b) \cap \mathbb{S}_{\infty} \ne \varnothing$, then $\mathbb{S}_{1}(b) \subset \mathbb{S}_{\infty}$.
\end{proposition}
\begin{proof}
    Suppose $f_{1}(\bm{a}) \in\mathbb{S}_{1}(b) \cap \mathbb{S}_{\infty}$.
    From \cref{cor:modulo-alpha-of-f_1}, $\forall k\in\mathbb{Z}$ such that $f_{1}(\bm{a} + k\Phi_1\bm{e}_1) > 0$, we have
    \[ f_{1}(\bm{a} + k\Phi_1\bm{e}_1) \equiv k\nabla_1 + f_{1}(\bm{a}) \equiv f_{1}(\bm{a})  \pmod{\alpha} \]
    Thus from \cref{prop:chr-are-periodic}, $f_{1}(\bm{a} + k\Phi_1\bm{e}_1) \in \mathbb{S}_{\infty}$. Notice that $f_{1}(\bm{a} + k\Phi_1\bm{e}_1)$ runs all $\mathbb{S}_{1}(b)$, thus $\mathbb{S}_{1}(b) \subset \mathbb{S}_{\infty}$.
\end{proof}

Actually, we can derive $\mathbb{V}_{n+1}(b)$ from $\mathbb{V}_{n}(b)$ and derive $\mathbb{S}_{n+1}(b)$ from $\mathbb{S}_{n}(b)$.
But since we are calculating the source value, the positive direction of the vector's dimensions is from the source value to the target value.
Therefore, if we want to obtain an $(n+1)$-dimensional vector from an $n$-dimensional vector, we should add the extra components to the beginning of the vector, rather than to the end.
Of course, the condition of $\nabla\cdot\mathfrak{V}_{n}$ should be satisfied anyway.
\begin{theorem}~\label{thm:layers-of-reduce-space-and-source-value-set}
    \begin{align*}
        \mathbb{S}_{n+1}(b) &= \bigcup_{a\in \mathbb{S}_{n}(b) - \mathbb{S}_{\infty}} \left\{\left. A_0 = \frac{\cor a  - \gamma}{\alpha}, A_{k+1} = \beta^{\Phi_1} A_{k} + \frac{\gamma}{\alpha}(\beta^{\Phi_1} - 1) \right\vert k\in\mathbb{N}\right\} \\
        &= \bigcup_{a\in \mathbb{S}_{n}(b) - \mathbb{S}_{\infty}} \left\{\left.\frac{1}{\alpha}\beta^{k\Phi_1}\cdot\cor a - \frac{\gamma}{\alpha}\right\vert k\in \mathbb{N}\right\},\\
        \mathbb{V}_{n+1}(b) &= \bigcup_{a\in \mathbb{S}_{n}(b) - \mathbb{S}_{\infty}} \left\{ (k\Phi_1 +  \chr a, \bm{v}) \mid  k\in \mathbb{N}, f_{n}(\bm{v};b) = a\right\}.
    \end{align*}
\end{theorem}
\begin{proof}
    If $\mathbb{S}_{n}(b) - \mathbb{S}_{\infty} = \varnothing$, then trivial.\newline
    Suppose $\mathbb{S}_{n}(b) - \mathbb{S}_{\infty} \ne \varnothing$. Then $\forall a\in \mathbb{S}_{n}(b) - \mathbb{S}_{\infty}$, we have $\chr a < \infty$. Then $A_0 = \frac{\cor a - \gamma}{\alpha} \in\widehat{\mathbb{N}}_{\beta}$, and hence $A_0 = f_{1}((\chr a);b)\in \mathbb{S}_{1}(a)\subset \mathbb{S}_{n+1}(b)$. From \cref{cor:single-vary-rule}, we have
    \begin{align*}
        A_{k+1} &= f_{1}(((k+1)\Phi_1 + \chr a);b) \\
        &= f_{1}((\Phi_1 + (k\Phi_1 + \chr a));b) \\
        &= \beta^{\Phi_1} A_{k} + \frac{\gamma}{\alpha}\cdot(\beta^{\Phi_1} - 1) \\
        &= \beta^{(k+1)\Phi_1} A_{0} + \frac{\gamma}{\alpha}\cdot(\beta^{(k+1)\Phi_1} - 1) \\
        &= \beta^{(k+1)\Phi_1}\cdot \frac{\cor a - \gamma}{\alpha} + \frac{\gamma}{\alpha}\cdot(\beta^{(k+1)\Phi_1} - 1) \\
        &= \beta^{(k+1)\Phi_1}\cdot \frac{\cor a}{\alpha} - \beta^{(k+1)\Phi_1}\cdot\frac{\gamma}{\alpha} + \frac{\gamma}{\alpha}\cdot\beta^{(k+1)\Phi_1} - \frac{\gamma}{\alpha} \\
        &= \frac{1}{\alpha}\cdot\beta^{(k+1)\Phi_1}\cdot\cor a - \frac{\gamma}{\alpha}.
    \end{align*}
    From \cref{thm:the-periodic-rule-of-reduce-sequence}, we have $A_{k+1}\in\mathbb{N}$, thus $A_{k+1}\in\mathbb{S}_{n+1}(b)$. \newline
    On the other hand, $\forall a\in\mathbb{S}_{n+1}(b)$, we have
    \[\beta^{\val(a;1) - \chr a_1}\cor a_1 = \alpha a + \gamma. \]
    We have $\val(a;1) - \chr a_1 \ge 0$. If not, then $\chr a_1$ is not the minimal valuation, or $\beta\mid(\beta^{\val(a;1)}a_1 - \gamma)/\alpha = a$. That's contradictory. So $\val(a;1) - \chr a_1 \ge 0$. And from \cref{thm:the-periodic-rule-of-reduce-sequence}, we have $\Phi_1\mid (\val(a;1) - \chr a_1)$. Denote $k = (\val(a;1) - \chr a_1)/\Phi_1$, since $a_1\in\mathbb{S}_{n}(b)$ by definition, then we have $a = \frac{1}{\alpha}\cdot\beta^{k\Phi_1}\cdot\cor a_1 - \frac{\gamma}{\alpha}$ in the set of the right side.\newline
    In conclusion, the theorem holds for $\mathbb{S}_{n + 1}(b)$. We can also verify from the above derivation that the conclusion regarding $\mathbb{V}_{n+1}(b)$ also holds.
\end{proof}
\begin{corollary}~\label{cor:source-value-set-equivlants-its-pre-ker-set}
    $\alpha\mathbb{S}_{n+1}(b) + \gamma = \mathbb{B}\otimes\cor\mathbb{S}_{n}(b)$. The decomposition is uniqe.
\end{corollary}
\begin{example}~\label{coro:layers-of-reduce-space-and-source-value-set-on-3-2-1}
    Suppose $(\alpha, \beta, \gamma) = (3,2,1)$.
    \begin{align*}
        \mathbb{S}_{n+1}(b) &= \bigcup_{a\in \mathbb{S}_{n}(b) - \mathbb{S}_{\infty}} \left\{A_0 = (\cor a  - 1)/3, A_{k+1} = 4f_{k} + 1 \mid k\in\mathbb{N}\right\} \\
        &= \bigcup_{a\in \mathbb{S}_{n}(b) - \mathbb{S}_{\infty}} \left\{\left.\frac{4^{k}}{3}\cdot\cor a - \frac{1}{3}\right\vert k\in \mathbb{N}\right\},\\
        \mathbb{V}_{n+1}(b) &= \bigcup_{a\in \mathbb{S}_{n+1}(b) - \mathbb{S}_{\infty}}\left\{ (2k + \bm{v}) \mid k\in \mathbb{N}, f_{n}(\bm{v};b) = a \right\}.
    \end{align*}
\end{example}
With this, we have
\begin{conjecture}~\label{conj:source-value-of-1-cover-all-odd-numbers}
    $\bigcup_{n=0}^{\infty}\mathbb{S}_n(1) = \lim_{n\to\infty}\mathbb{S}_n(1) = \widehat{\mathbb{N}}_{{2}}$.
\end{conjecture}
We can verify it equivalents to the classic Collatz problem.

\cref{thm:layers-of-reduce-space-and-source-value-set} provides a very important insight: the source values of a number are not entirely chaotic; they can be precisely categorized into different levels.
However, although we have provided expressions for $\mathbb{V}_{n+1}(b)$ and $\mathbb{S}_{n+1}(b)$, they are still complex and difficult to use.
They are merely another equivalent expression of the original problem and are hard to use for finding counter examples.
That's because they are still based on $\mathbb{S}_{n}(b)$, and therefore still require recursively representing all the first $n$ $\mathbb{S}$s to represent the $(n+1)$-th. To solve this problem, we will try to decompose the reduce vectors in \cref{sec:decomposition-of-reduce-vector} and provide an other view on them source value set.

We can observe that as the dimension $n$ increases, $\mathbb{V}_{n}(b)$ becomes increasingly "sparse," as the terminal components—those lying between two integers—increment in multiples of $\Phi_{n}$. However, this represents merely a local perspective focused on a single value within a single hierarchical level. In reality—given that numerous numbers have already been situated within lower-dimensional spaces, and that the original source itself constrains the expansion of dimensions—the actual number of elements within the set $\mathbb{S}_{n}(b)$ undergoes an exponential expansion; this phenomenon is dictated by the recursive construction described in \cref{thm:layers-of-reduce-space-and-source-value-set}, which exactly is because it's "complex and difficult to use".

Looking back at \cref{thm:the-periodic-rule-of-reduce-sequence} from \cref{thm:layers-of-reduce-space-and-source-value-set}, we can finally understand that when $\Phi_k$ adjusts the reduction sequence, we are actually changing one number in $\mathbb{S}_{n-k}(b)$ to another, while keeping the beginning and subsequent parts of the reduction sequence unchanged. It also explains the role of $\gamma$. It shows that $\gamma$ doesn't participate much in the construction of the source value set, but it determines the size of the jump steps. Although we might find it seems that only $\Phi_1$ participates in the construction of the source value set, in reality, all $\Phi(k)$ are power of $\alpha$ multiplied by $\Phi(1)$. Therefore, $\Phi_k$ not only represents the construction of the source value set but also tells us about the cyclic structure of the source value set in the form of factors of $\alpha$.
Note that both numbers also have their own source value sets, which implies some homomorphisms between these two source value sets.
And the two theorems reveal that elements of $\mathbb{S}_{1}(b)$ is in a same class and their $\val$ are linear recurrenced.

\begin{theorem}~\label{thm:infinite-layers}
	Suppose $\alpha^{2}\mid (\beta^{\Phi_1} - 1)$. For $b\not\in\mathbb{S}_{\infty}$, we have $\forall N > 0$, there is $\mathbb{S}(b) - \bigcup_{n=0}^{N}\mathbb{S}_{n}(b) \ne \varnothing$.
\end{theorem}
\begin{proof}
	Divide into two situations:\newline
	If $b \not\rightsquigarrow b$. Since $b\not\in\mathbb{S}_{\infty}$, from \cref{prop:gcd-1-derive-not-all-primitive}, we have $\mathbb{S}_{N}(b) \not\subset \mathbb{S}_{\infty}$. That means $\mathbb{S}_{N + 1}(b) \ne \varnothing$.
	From \cref{prop:non-cycle-layer-independent}, $\forall N \ge 0$ there is
	\[\mathbb{S}(b) - \bigcup_{n=0}^{N}\mathbb{S}_{n}(b) \supset \mathbb{S}_{N+1}(b) - \bigcup_{n=0}^{N}\mathbb{S}_{n}(b) = \mathbb{S}_{N+1}(b) \ne \varnothing.\]
	If $b \rightsquigarrow b$. Suppose it's a $t$-period cycle, then $b \in\mathbb{S}_{t}(b)$. From \cref{thm:layers-of-reduce-space-and-source-value-set}, we can just pick $a = \beta^{\Phi_1}\cdot b + \frac{\gamma}{\alpha}\cdot(\beta^{\Phi_1} - 1)\in\mathbb{S}_{t}(b)$, we have $a \rightsquigarrow b \not\rightsquigarrow a$. From previous discussion, $\mathbb{S}(a)$ can not be covered within finity layers, thus the $\mathbb{S}(b) \supset \mathbb{S}(a)$ also can't be covered within finity layers.
\end{proof}
\cref{thm:infinite-layers} demonstrates that we can never resolve the Collatz conjecture through any method of case analysis by number's patterns, even if we partition the problem into various infinite sets. It's because different layers are in defferent patterns. Though we can tell a lot of patterns, and for any given layer, we can also explore its pattern, that's not enough at all.

Finally, let us introduce two notations to provide a perspective on a single source value and reinterpret the theorem \cref{thm:layers-of-reduce-space-and-source-value-set}.
\begin{definition}
    $S(k; b) = \frac{1}{\alpha}\beta^{k\Phi_1}\cdot\cor b - \frac{\gamma}{\alpha}$, which is the $k$-th element of $\mathbb{S}_1(b)$.
\end{definition}
\begin{proposition}
    $S(k + 1; b) > S(k; b)$.
\end{proposition}
\begin{definition}
    Define $F: \mathbb{S}_n(b)\to\mathbb{S}_{n+1}(b)$ that $F(a) = S(0; a)$ and define $\Phi: \mathbb{S}_1(a)\to \mathbb{S}_1(a)$ that $\Phi(S(m;a)) = S(m + 1;a)$.
\end{definition}
\begin{theorem}~\label{thm:f_n-by-F-and-Phi}
    $\forall f_{n}(\bm{a};b)\in\mathbb{S}(b)$, we have $f_{n}(\bm{a};b) = \Phi^{k_{1}}F\Phi^{k_{2}}\cdots F\Phi^{k_{n}}F(b)$.
\end{theorem}
